\section{Journal Guides \& Positions}

\subsection{6 News Australia}

This outlet is a digitally native outlet in Australia, which was founded in 2018. They built their audience mostly through social media, catering to younger Australians, especially with coverage of elections and politics.\footnote{6 News Australia. (n.d.). About. \url{https://www.6newsaustralia.com/about}} The outlet is centrist and non-partisan, and prioritises factual reporting over opinion commentary.\footnote{Reuters Institute for the Study of Journalism. (2025). Digital News Report 2025. University of Oxford. \url{https://reutersinstitute.politics.ox.ac.uk/digital-news-report/2025}}

\subsection{Al-Ahram}

This outlet is Egypt’s oldest newspaper, which was founded in 1875. The outlet is state-owned, which makes it a semi-official voice of the Egyptian government on both domestic and international affairs.\footnote{Rugh, W. A. (2004). Arab Mass Media: Newspapers, Radio, and Television in Arab Politics. Praeger} Their coverage focuses on politics in the Middle East, African and Arab affairs. Their political position is often nationalist, pro-sovereignty, and non-intervention.\footnote{Reporters Without Borders. (2025). Egypt. \url{https://rsf.org/en/country/egypt}}

\subsection{Al Jazeera}

This outlet is a Qatari state-funded international news network. It was launched in 1996. Now, it has amassed an audience of over 270 million households, and 80 offices worldwide.\footnote{Al Jazeera Media Network. (n.d.). About us. \url{https://www.aljazeera.com/about-us}} Its political position is often associated with the Qatari foreign policy interests, but is also claimed to be slightly left on the general political spectrum.\footnote{AllSides. (n.d.). Al Jazeera media bias rating. \url{https://www.allsides.com/news-source/al-jazeera-media-bias}}

\subsection{Asahi Shimbun}

This outlet is one of Japan’s five major daily news organisations. It was founded in 1879, and now has one of the highest circulations in the world.\footnote{Asahi Shimbun. (n.d.). Corporate profile. \url{https://www.asahi.com/ajw/corporate/}} By Japanese standards, it is considered centre-left politically, and advocates for civil liberties, pacifism and constitutional democracy.\footnote{Huffman, J. L. (1997). Creating a Public: People and Press in Meiji Japan. University of Hawaii Press}

\subsection{BBC News}

This outlet is the news division of the UK’s national broadcaster, which is publicly funded. Founded in 1922, it began with a Royal Charter requiring impartiality, accuracy, and balance.\footnote{BBC. (n.d.). Inside the BBC: Mission, values and public purposes. \url{https://www.bbc.com/aboutthebbc}} Its stance is centrist, and internally reflects a liberal-democratic worldview, focusing on human rights and the rule of law.\footnote{AllSides. (2024). How Americans rated the bias of BBC, Reuters, CNN, Fox News, and New York Post. \url{https://www.allsides.com/blind-survey/how-americans-rated-bias-bbc-reuters-cnn-fox-news-and-new-y} ork-post}

\subsection{Corriere Della Sera}

This outlet is Italy’s most ready daily newspaper. It was founded in Milan in 1876.\footnote{Corriere della Sera. (n.d.). Chi siamo. \url{https://www.corriere.it/chi-siamo/}} The outlet is considered the country’s newspaper of record, and has a centre to centre-right position. It favours pro-European integration and liberal economics.\footnote{Reuters Institute for the Study of Journalism. (2025). Digital News Report 2025: Italy. University of Oxford.}

\subsection{Deutsche Welle}

This outlet is Germany’s international broadcaster, which is publicly funded. It was established in 1953, and now broadcasts in 32 languages.\footnote{Deutsche Welle. (n.d.). About DW. \url{https://www.dw.com/en/about-dw/s-30600}} It has a centrist, pro-European and pro-democratic stance and emphasises pieces on human rights, the rule of law, and promotes multilateral democracy.\footnote{Deutsche Welle Act. (2004). \url{https://www.dw.com/en/dw-act/a-18265246}}

\subsection{El Pais}

This outlet is Spain’s highest-circulation daily newspaper. It was founded in 1976 during Spain’s transition to democracy. It is widely influential across Spanish-speaking countries, largely including Latin America.\footnote{El País. (n.d.). Quiénes somos. \url{https://elpais.com/corporativos/quienes-somos/}} Its position is centre-left and is often aligned with social democracy, European federalism, and internationalism. \footnote{Reuters Institute for the Study of Journalism. (2025). Digital News Report 2025: Spain. University of Oxford.}

\subsection{Fox News}

This outlet is the most watched news network on American cable television. It was launched in 1996, and has a nationalist, right-wing, conservative stance.\footnote{AllSides. (n.d.). Fox News media bias rating. \url{https://www.allsides.com/news-source/fox-news-media-bias}} Its international reporting often focuses on American interests, especially a strong stance on defence and being sceptical of multilateral diplomacy.\footnote{Grynbaum, M. M. (2023, April 18). Dominion and Fox News settle defamation case for \$787.5 million. The New York Times. \url{https://www.nytimes.com/2023/04/18/business/fox-news-dominion-settlement.html}}

\subsection{Le Monde}

This outlet is France’s newspaper of record, and was founded in 1944. It has a centre-left stance and a strong culture of investigative journalism.\footnote{Le Monde. (n.d.). Qui sommes-nous? \url{https://www.lemonde.fr/qui-sommes-nous/}} It promotes multilateralism, international law and collaborative European strategy. It is sceptical of American hegemony and Russian aggression.\footnote{Reuters Institute for the Study of Journalism. (2025). Digital News Report 2025: France. University of Oxford.}

\subsection{Nunatsiaq}

This outlet is an independent newspaper which caters to Canada's Arctic and Inuit communities. It was founded in 1973, and is the primary English-language news organisation for the territory.\footnote{Nunatsiaq News. (n.d.). About us. \url{https://nunatsiaq.com/about/}} It does not have a partisan position, but it does center on Indigenous perspectives, and covers topics such as Arctic sovereignty, climate change, and how policies affect northern communities.\footnote{Basen, I. (2019). Covering the Arctic: A guide to Northern journalism. Canadian Journalism Foundation.}

\subsection{O Globo}

This outlet is one of Brazil’s largest daily newspapers. It was founded in 1925, and is part of Grupo Globo (Latin America’s largest media conglomerate).\footnote{Grupo Globo. (n.d.). Who we are. \url{https://grupoglobo.globo.com/en/who-we-are/}} Their position is centre-right, and favours market-oriented economics. It reflects Brazil’s leading role for non-Western countries in diplomatic groups such as BRICS and the G20.\footnote{Reuters Institute for the Study of Journalism. (2025). Digital News Report 2025: Brazil. University of Oxford.}

\subsection{Reuters}

This outlet is one of the world’s oldest and largest news agencies. It was founded in 1851, and is widely regarded as a benchmark for impartial reporting. Its political stance is largely centrist and focuses on a framework of values such as independence, integrity, and freedom from bias.\footnote{AllSides. (2024). How Americans rated the bias of BBC, Reuters, CNN, Fox News, and New York Post; Thomson Reuters. (n.d.). Trust Principles. \url{https://www.thomsonreuters.com/en/about-us/trust-principles.html}}

\subsection{South China Morning Post}

This outlet is based in Hong Kong, and is the leading English-language newspaper for the area. It was founded in 1903 and has been owned since 2016 by Alibaba Group.\footnote{South China Morning Post. (n.d.). About us. \url{https://www.scmp.com/about}} It claims to focus on professional journalistic standards, but there have been concerns due to Hong Kong’s political environment and Chinese influence on reporting of sensitive topics.\footnote{Reporters Without Borders. (2025). Hong Kong. \url{https://rsf.org/en/country/hong-kong}; Cook, S. (2020). Beijing's Global Megaphone: The Expansion of Chinese Communist Party Media Influence since 2017. Freedom House.}

\subsection{TASS}

This outlet is Russia’s largest state-owned news agency. It was founded in 1904, and is the primary vehicle for Kremlin communication to both domestic and international audiences.\footnote{TASS. (n.d.). About TASS. \url{https://tass.com/about}} Its political stance closely follows Russian state interests, and emphasises sovereignty, while criticising Western states and their sanctions.\footnote{Reporters Without Borders. (2025). Russia. \url{https://rsf.org/en/country/russia}}

\subsection{The Australian}

This outlet is a national broadsheet, which was founded in 1964. It was founded by Rupert Murdoch, and is a part of News Corp Australia. Its stance is centre-right with a strong position on economic liberalism and Australian national security.\footnote{The Australian. (n.d.). About The Australian. \url{https://www.theaustralian.com.au/about}}

\subsection{The Daily Nation}

This outlet is Kenya’s largest newspaper. It was founded in 1960 and is published by the Nation Media Group, which is the largest media conglomerate in East Africa.\footnote{Nation Media Group. (n.d.). About us. \url{https://www.nationmedia.com/about-us/}} Its position is centrist and mostly independent. It often writes about the tension between diplomacy by global powers and respect for sovereignty and development.\footnote{Reporters Without Borders. (2025). Kenya. \url{https://rsf.org/en/country/kenya}; Reuters Institute for the Study of Journalism. (2025). Digital News Report 2025: Kenya. University of Oxford.}

\subsection{The Economist}

This outlet is a weekly magazine published in Britain. It was founded in 1843, and is globally recognised as a leader on reporting international affairs, economics, and politics.\footnote{The Economist. (n.d.). About us. \url{https://www.economist.com/about-the-economist}} Its stance is classically liberal, and focuses on the free market, globalisation, and a rules-based international order. It supports NATO and the EU. However, it often states it wishes to maintain its intellectual independence from states.\footnote{The Economist. (2018, September 13). What The Economist believes. \url{https://www.economist.com/leaders/2018/09/13/what-the-economist-believes}}

\subsection{The Times of India}

This outlet is the world’s largest selling English language daily newspaper. It was founded in 1838, and its position is centrist. It often reflects on India’s position as a major non-aligned power within global politics, as it maintains a relationship with the West and Russia.\footnote{The Times of India. (n.d.). About us. \url{https://timesofindia.indiatimes.com/aboutus.cms}}

\subsection{Vox}

This outlet is an American digital media company, which was founded in 2014. It is most known for explanatory journalism, attempting to make complex policy decisions accessible to a general audience.\footnote{Vox Media. (n.d.). About Vox. \url{https://www.vox.com/pages/about}} Its stance is centre-left or sometimes considered progressive. It emphasises social justice, climate policy and evidence-base analysis.\footnote{AllSides. (n.d.). Vox media bias rating. \url{https://www.allsides.com/news-source/vox-media-bias}}
